\setlength\cwii{\b}
\setlength\cwiii{\c}
\setlength\cwi{\a}
\begin{tcbitemize}[raster columns=3,raster rows=7,raster equal height=rows,
            raster force size=false,
            raster column skip=-0.8pt,raster row skip=-0.8pt,
        raster column 1/.style={width=\cwi,
            fontupper=\bfseries},
        raster column 2/.style={width=\cwii},
        raster column 3/.style={width=\cwiii,fontupper=\bfseries,
            halign=center,valign=center},
        ]
    \tcbitem 1)
    \tcbitem 
             \begin{enumerate}[label=\textbf{(\arabic*)}]
                 \item pH-mètre
                 \item Solution d'hydroxyde de sodium
                       (solution titrante)
                 \item Solution d'acide azélaïque
                       (solution titrée)
             \end{enumerate}
    \tcbitem (0,75pt)
    \tcbitem 2)
    \tcbitem $\ce{AH_{(aq)} + HO^-_{(aq)} -> A^-_{(aq)} + H2O_{(\ell)}}$
    \tcbitem (0,5pt)
    \tcbitem 3)
    \tcbitem $V_{B\mathrm{E}}=\qty{10.0}{\mL} \quad;\quad \mathrm{pH_{E}}=8,1$
    \tcbitem (1pt)
    \tcbitem 4)
    \tcbitem $C\times V_A=C_B\times V_{B\mathrm{E}}\quad\Rightarrow\;
             C=\frac{C_B\times V_{B\mathrm{E}}}{V_A}=
             \frac{\num{1,00e-2}\times\num{10,0}}{\num{10,0}}=
             \qty{1.00e-2}{\M}$
    \tcbitem (1pt)
    \tcbitem 5)
    \tcbitem[raster multicolumn=2,raster multirow=4,blankest]
        \begin{tcbitemize}[raster rows=4,raster columns=2,
               raster height=\tcbtextheight,
               raster column 1/.style={width=\cwii},
               raster column 2/.style={width=\cwiii,fontupper=\bfseries,
                valign=center}
            ]
            \tcbitem L'indicateur convenable est le \textbf{rouge de crésol}.
            \tcbitem (0,25pt)
            \tcbitem \textbf{Justification~:}
                     sa zone de virage encadre le $\mathrm{pH}$ à l'équivalence
                     ($8,1$).
            \tcbitem (0,25pt)
        \end{tcbitemize}
\end{tcbitemize}
